% GENERATED FILE. Edit canonical lessons, then run npm run generate:books.
% canonical-source-hash: fnv1a64:4e290102d1413704
% canonical-lessons: ES-C718-guinar, ES-C718-bromista, ES-C718-broma, ES-C718-bromear, ES-C718-burlarse

\chapter{Guiños y bromas}
\label{ch:es-guinos-y-bromas}
\hlchaptermodality{718}

\begin{chapteropening}
\textbf{By the end of this chapter:} \emph{I can talk about winking, joking and teasing among friends.}

Complete the last lesson of chapter 718: I can talk about winking, joking and teasing among friends.
\end{chapteropening}

\section[guiñar]{guiñar — to wink}

\label{lesson:ES-C718-guinar}



\begin{quote}
Before the new one: say the Spanish for the tear, then the Spanish for desperate.
\end{quote}

\subsection*{You'll want to know: guiñar}
\textbf{guiñar} — \textquotedblleft{}to wink\textquotedblright{}.

\begin{quote}
Me guiñó un ojo y sonreí.
\end{quote}

\begin{quote}
\emph{He winked at me and I smiled.}
\end{quote}

\emph{Guiñar un ojo} — to wink. \emph{Un guiño} — a wink.

\begin{cousinweb}
For \emph{guiñar} the origin is not settled; a Germanic source has been suggested, and some see an imitative word, but none is proved.
\end{cousinweb}

\begin{grammarlens}[title={What you've built}]
One more word for telling, explaining and describing at B1.
\end{grammarlens}

\subsection*{Your turn}
\begin{itemize}
\raggedright
  \item \emph{Say it:} \emph{guiñar}
  \item \emph{Say it:} \emph{guiñar}, once more
  \item \emph{Say it:} a sentence of your own with \emph{guiñar}
  \item \emph{Recall:} say the Spanish for the tear, then the Spanish for desperate, then say \emph{guiñar} again
\end{itemize}

\begin{tcolorbox}[breakable,colback=teal!4,colframe=teal!35!black,title={Before you move on}]
What does guiñar mean? (\textquotedblleft{}To wink\textquotedblright{}.) Say it once more.
\end{tcolorbox}
\section[bromista]{bromista — joker; fond of jokes}

\label{lesson:ES-C718-bromista}



\begin{quote}
Before the new one: say the Spanish for desperate, then the Spanish for to wink.
\end{quote}

\subsection*{You'll want to know: bromista}
\textbf{bromista} — \textquotedblleft{}joker; fond of jokes\textquotedblright{}.

\begin{quote}
Mi tío es muy bromista.
\end{quote}

\begin{quote}
\emph{My uncle loves a joke.}
\end{quote}

Same form for men and women: \emph{es una bromista}.

\begin{cousinweb}
\emph{bromista} is built on \emph{broma} \textquotedblleft{}joke\textquotedblright{} with the ending \emph{-ista}; \emph{broma} comes from Greek \emph{brôma} \textquotedblleft{}food\textquotedblright{} (EtymDB).
\end{cousinweb}

\begin{grammarlens}[title={What you've built}]
One more word for telling, explaining and describing at B1.
\end{grammarlens}

\subsection*{Your turn}
\begin{itemize}
\raggedright
  \item \emph{Say it:} \emph{bromista}
  \item \emph{Say it:} \emph{bromista}, once more
  \item \emph{Say it:} a sentence of your own with \emph{bromista}
  \item \emph{Recall:} say the Spanish for desperate, then the Spanish for to wink, then say \emph{bromista} again
\end{itemize}

\begin{tcolorbox}[breakable,colback=teal!4,colframe=teal!35!black,title={Before you move on}]
What does bromista mean? (\textquotedblleft{}Joker; fond of jokes\textquotedblright{}.) Say it once more.
\end{tcolorbox}
\section[la broma]{la broma — the joke, the prank}

\label{lesson:ES-C718-broma}



\begin{quote}
Before the new one: say the Spanish for to wink, then the Spanish for joker.
\end{quote}

\subsection*{You'll want to know: la broma}
\textbf{la broma} — \textquotedblleft{}the joke, the prank\textquotedblright{}.

\begin{quote}
Era una broma, no te enfades.
\end{quote}

\begin{quote}
\emph{It was a joke, do not be angry.}
\end{quote}

\emph{Gastar una broma} — to play a joke. \emph{En broma} — jokingly.

\begin{cousinweb}
\emph{broma} comes from Greek \emph{brôma} \textquotedblleft{}food\textquotedblright{} (EtymDB), through Latin. It first named the shipworm that eats ship timber, then any nuisance, and at last a joke.
\end{cousinweb}

\begin{grammarlens}[title={What you've built}]
One more word for telling, explaining and describing at B1.
\end{grammarlens}

\subsection*{Your turn}
\begin{itemize}
\raggedright
  \item \emph{Say it:} \emph{la broma}
  \item \emph{Say it:} \emph{la broma}, once more
  \item \emph{Say it:} a sentence of your own with \emph{la broma}
  \item \emph{Recall:} say the Spanish for to wink, then the Spanish for joker, then say \emph{la broma} again
\end{itemize}

\begin{tcolorbox}[breakable,colback=teal!4,colframe=teal!35!black,title={Before you move on}]
What does la broma mean? (\textquotedblleft{}The joke, the prank\textquotedblright{}.) Say it once more.
\end{tcolorbox}
\section[bromear]{bromear — to joke}

\label{lesson:ES-C718-bromear}



\begin{quote}
Before the new one: say the Spanish for joker, then the Spanish for the joke.
\end{quote}

\subsection*{You'll want to know: bromear}
\textbf{bromear} — \textquotedblleft{}to joke\textquotedblright{}.

\begin{quote}
No te preocupes, estoy bromeando.
\end{quote}

\begin{quote}
\emph{Do not worry, I am joking.}
\end{quote}

\emph{Estar bromeando} — to be joking.

\begin{cousinweb}
\emph{bromear} is built on \emph{broma} \textquotedblleft{}joke\textquotedblright{} with the verb ending \emph{-ear}; \emph{broma} comes from Greek \emph{brôma} \textquotedblleft{}food\textquotedblright{} (EtymDB).
\end{cousinweb}

\begin{grammarlens}[title={What you've built}]
One more word for telling, explaining and describing at B1.
\end{grammarlens}

\subsection*{Your turn}
\begin{itemize}
\raggedright
  \item \emph{Say it:} \emph{bromear}
  \item \emph{Say it:} \emph{bromear}, once more
  \item \emph{Say it:} a sentence of your own with \emph{bromear}
  \item \emph{Recall:} say the Spanish for joker, then the Spanish for the joke, then say \emph{bromear} again
\end{itemize}

\begin{tcolorbox}[breakable,colback=teal!4,colframe=teal!35!black,title={Before you move on}]
What does bromear mean? (\textquotedblleft{}To joke\textquotedblright{}.) Say it once more.
\end{tcolorbox}
\section[burlarse]{burlarse — to make fun of}

\label{lesson:ES-C718-burlarse}



\begin{quote}
Before the new one: say the Spanish for the joke, then the Spanish for to joke.
\end{quote}

\subsection*{You'll want to know: burlarse}
\textbf{burlarse} — \textquotedblleft{}to make fun of\textquotedblright{}.

\begin{quote}
No te burles de tu hermano.
\end{quote}

\begin{quote}
\emph{Do not make fun of your brother.}
\end{quote}

\emph{Burlarse de alguien}. Less kind than \emph{bromear}.

\begin{cousinweb}
\emph{burlarse} is built on \emph{burla} \textquotedblleft{}mockery\textquotedblright{}. For \emph{burla} the origin is not settled; one old guess links it with Latin \emph{burrae} \textquotedblleft{}trifles\textquotedblright{}.
\end{cousinweb}

\begin{grammarlens}[title={What you've built}]
One more word for telling, explaining and describing at B1.
\end{grammarlens}

\subsection*{Your turn}
\begin{itemize}
\raggedright
  \item \emph{Say it:} \emph{burlarse}
  \item \emph{Say it:} \emph{burlarse}, once more
  \item \emph{Say it:} a sentence of your own with \emph{burlarse}
  \item \emph{Recall:} say the Spanish for the joke, then the Spanish for to joke, then say \emph{burlarse} again
\end{itemize}

\begin{tcolorbox}[breakable,colback=teal!4,colframe=teal!35!black,title={Before you move on}]
What does burlarse mean? (\textquotedblleft{}To make fun of\textquotedblright{}.) Say it once more.
\end{tcolorbox}
